\section{Задание}
Построить для данной табличной функции интерполяционный кубический сплайн дефекта 1 с заданными граничными условиями. Коэффициенты вычислять с четырьмя знаками после запятой в мантиссе. Проверить все условия, налагаемые на сплайн.

\section{Цель работы}
Научиться строить интерполяционные кубические сплайны для табличных функций и проверять условия, налагаемые на построенный сплайн.

\section{Ход работы}

\textit{Перед выполнением работы отметим, что все исходные данные взяты из варианта 47.}

Для заданной табличной функции требуется построить интерполяционный кубический сплайн дефекта 1. Значения функции приведены в таблице \ref{table:values}.

\begin{table}[H]
    \centering
    \caption{Значения $x_i,y_i$ для построения кубического сплайна}
    \label{table:values}
    \begin{tblr}{
        colspec = {X[-1,l]X[c]X[c]X[c]X[c]X[c]X[c]X[c]X[c]X[c]|},
        hlines,
        vlines,
        rows = {abovesep=8pt,belowsep=8pt,font=\small $},
        column{1} = {leftsep=8pt},
        }
        x_i & 0.050 & 0.052 & 0.060 & 0.065 & 0.069 & 0.075 & 0.085 & 0.090 & 0.096 \\
        y_i & 0.2079 & 0.2081 & 0.2090 & 0.2095 & 0.2099 & 0.2105 & 0.2116 & 0.2121 & 0.2127 \\
    \end{tblr}
\end{table}

По условию функция считается периодической с периодом
\[
T=x_n-x_0=0.096-0.050=0.046.
\]

Обозначим шаги таблицы
\[
h_i=x_{i+1}-x_i, \quad i=0,1,\dots,n-1,
\]

где $n + 1 = 9 \iff n = 8.$

Для данной таблицы получим
\[
\begin{aligned}
&h_0=0.002, \quad h_1=0.008, \quad h_2=0.005, \quad h_3=0.004,\\
&h_4=0.006, \quad h_5=0.010, \quad h_6=0.005, \quad h_7=0.006.
\end{aligned}
\]

\subsection{Краткая теория}

Кубический интерполяционный сплайн строится как набор кубических полиномов
\[
S_i(x), \quad x \in [x_i, x_{i+1}], \quad i = 0, 1, \dots, n-1,
\]
которые удовлетворяют следующим условиям:

\begin{itemize}
    \item условия интерполяции:
    \[
    S_i(x_i) = y_i, \quad S_i(x_{i+1}) = y_{i+1};
    \]

    \item непрерывность первой и второй производных:
    \[
    S'_i(x_{i+1}) = S'_{i+1}(x_{i+1}), \quad
    S''_i(x_{i+1}) = S''_{i+1}(x_{i+1});
    \]

    \item граничные условия (в данной работе — периодические):
    \[
    S'(x_0) = S'(x_n), \quad
    S''(x_0) = S''(x_n).
    \]
\end{itemize}

Каждый частичный полином имеет вид:
\[
S_i(x) = a_i + b_i(x - x_i) + c_i(x - x_i)^2 + d_i(x - x_i)^3.
\]

В стандартной форме используются значения второй производной в узлах:
\[
M_i = S''(x_i).
\]

Для них получается система уравнений:
\[
h_{i-1}M_{i-1} + 2(h_{i-1} + h_i)M_i + h_i M_{i+1}
=
6\left(
\frac{y_{i+1} - y_i}{h_i}
-
\frac{y_i - y_{i-1}}{h_{i-1}}
\right),
\]
\[
i = 1, \dots, n-1,
\]

где
\[
h_i = x_{i+1} - x_i.
\]

С учётом периодических граничных условий:
\[
M_0 = M_n.
\]

После нахождения $M_i$ коэффициенты сплайна вычисляются по формулам:
\[
a_i = y_i,
\]
\[
b_i =
\frac{y_{i+1} - y_i}{h_i}
- \frac{h_i}{6}(2M_i + M_{i+1}),
\]
\[
c_i = \frac{M_i}{2},
\quad
d_i = \frac{M_{i+1} - M_i}{6h_i}.
\]

\subsection{Реализация}

На каждом частичном отрезке $[x_i;x_{i+1}]$ будем искать кубический сплайн в виде
\[
S_i(x)=a_i+b_i(x-x_i)+c_i(x-x_i)^2+d_i(x-x_i)^3,
\quad x\in[x_i;x_{i+1}].
\]

Коэффициенты выражаются через значения функции и наклоны $m_i=S'(x_i)$:
\[
a_i=y_i, \quad b_i=m_i,
\]
\[
c_i=\frac{3\Delta_i-2m_i-m_{i+1}}{h_i},
\quad
d_i=\frac{m_i+m_{i+1}-2\Delta_i}{h_i^2},
\]
где
\[
\Delta_i=\frac{y_{i+1}-y_i}{h_i}.
\]

Найдём значения $\Delta_i$:
\[
\begin{aligned}
&\Delta_0=0.1000, \quad \Delta_1=0.1125, \quad \Delta_2=0.1000, \quad \Delta_3=0.1000,\\
&\Delta_4=0.1000, \quad \Delta_5=0.1100, \quad \Delta_6=0.1000, \quad \Delta_7=0.1000.
\end{aligned}
\]

Для внутренних узлов условия непрерывности второй производной имеют вид
\[
h_i m_{i-1}+2(h_{i-1}+h_i)m_i+h_{i-1}m_{i+1}
=3\left(\frac{h_i}{h_{i-1}}(y_i-y_{i-1})+\frac{h_{i-1}}{h_i}(y_{i+1}-y_i)\right),
\]
где $i=1,2,\dots,n-1$.

Так как функция считается периодической, добавим граничные условия
\[
m_0=m_n,
\]
\[
S''_0(x_0)=S''_{n-1}(x_n).
\]

В результате получаем систему линейных уравнений относительно неизвестных наклонов $m_0,m_1,\dots,m_8$:
\[
\left(
\begin{array}{ccccccccc}
0.008 & 0.020 & 0.002 & 0 & 0 & 0 & 0 & 0 & 0 \\
0 & 0.005 & 0.026 & 0.008 & 0 & 0 & 0 & 0 & 0 \\
0 & 0 & 0.004 & 0.018 & 0.005 & 0 & 0 & 0 & 0 \\
0 & 0 & 0 & 0.006 & 0.020 & 0.004 & 0 & 0 & 0 \\
0 & 0 & 0 & 0 & 0.010 & 0.032 & 0.006 & 0 & 0 \\
0 & 0 & 0 & 0 & 0 & 0.005 & 0.030 & 0.010 & 0 \\
0 & 0 & 0 & 0 & 0 & 0 & 0.006 & 0.022 & 0.005 \\
1 & 0 & 0 & 0 & 0 & 0 & 0 & 0 & -1 \\
-1000 & -500 & 0 & 0 & 0 & 0 & 0 & -166.6667 & -333.3333
\end{array}
\right)
\left(
\begin{array}{c}
m_0\\m_1\\m_2\\m_3\\m_4\\m_5\\m_6\\m_7\\m_8
\end{array}
\right)
=
\left(
\begin{array}{c}
0.0031\\0.0041\\0.0027\\0.0030\\0.0050\\0.0047\\0.0033\\0\\-200.0000
\end{array}
\right).
\]

На листинге \ref{lst:spline} приведена программа, вычисляющая коэффициенты кубического сплайна.

\begin{listing}[H]
\caption{Вычисление коэффициентов периодического кубического сплайна}
\vspace{-20pt}
\label{lst:spline}
\begin{minted}[linenos, numbersep=10pt, firstnumber=1]{matlab}
x = [0.050, 0.052, 0.060, 0.065, 0.069, 0.075, 0.085, 0.090, 0.096];
y = [0.2079, 0.2081, 0.2090, 0.2095, 0.2099, 0.2105, 0.2116, 0.2121, 0.2127];

h = diff(x);
delta = diff(y) ./ h;
n = length(x) - 1;

A = zeros(n + 1, n + 1);
b = zeros(n + 1, 1);

for i = 2:n
    A(i - 1, i - 1) = h(i);
    A(i - 1, i) = 2 * (h(i - 1) + h(i));
    A(i - 1, i + 1) = h(i - 1);
    b(i - 1) = 3 * (h(i) / h(i - 1) * (y(i) - y(i - 1)) + ...
        h(i - 1) / h(i) * (y(i + 1) - y(i)));
end

% Periodic condition: S'(x0) = S'(xn)
A(n, 1) = 1;
A(n, n + 1) = -1;
b(n) = 0;

% Periodic condition: S''(x0) = S''(xn)
A(n + 1, 1) = -2 / h(1);
A(n + 1, 2) = -1 / h(1);
A(n + 1, n) = -1 / h(n);
A(n + 1, n + 1) = -2 / h(n);
b(n + 1) = -3 * delta(1) / h(1) - 3 * delta(n) / h(n);

m = A \ b;

for i = 1:n
    a(i) = y(i);
    bcoef(i) = m(i);
    c(i) = (3 * delta(i) - 2 * m(i) - m(i + 1)) / h(i);
    d(i) = (m(i) + m(i + 1) - 2 * delta(i)) / h(i)^2;
end

disp([a' bcoef' c' d'])
\end{minted}
\end{listing}

В результате решения системы были получены следующие значения наклонов:
\[
\begin{aligned}
&m_0=0.0988, \quad m_1=0.1035, \quad m_2=0.1070, \quad m_3=0.0986,\\
&m_4=0.0994, \quad m_5=0.1050, \quad m_6=0.1045, \quad m_7=0.0991, \quad m_8=0.0988.
\end{aligned}
\]

Следовательно, коэффициенты кубического сплайна имеют вид, представленный в таблице \ref{table:coeffs}.

\begin{table}[H]
    \centering
    \caption{Коэффициенты кубического сплайна}
    \label{table:coeffs}
    \begin{tblr}{
        colspec = {X[c]X[c]X[c]X[c]X[c]},
        hlines,
        vlines,
        rows = {abovesep=5pt,belowsep=5pt, $},
        row{1} = {font=\normalfont}
    }
        i & a_i & b_i & c_i & d_i \\
        0 & 0.2079 & 0.0988 & -0.5609 & 582.3924 \\
        1 & 0.2081 & 0.1035 & 2.9335 & -226.6399 \\
        2 & 0.2090 & 0.1070 & -2.5059 & 222.8726 \\
        3 & 0.2095 & 0.0986 & 0.8372 & -122.6772 \\
        4 & 0.2099 & 0.0994 & -0.6349 & 121.8499 \\
        5 & 0.2105 & 0.1050 & 1.5584 & -105.4709 \\
        6 & 0.2116 & 0.1045 & -1.6058 & 141.5796 \\
        7 & 0.2121 & 0.0991 & 0.5179 & -59.9346 \\
    \end{tblr}
\end{table}

Таким образом, интерполяционный кубический сплайн имеет вид
\[
S(x)=
\begin{cases}
S_0(x), & x\in[0.050;0.052],\\
S_1(x), & x\in[0.052;0.060],\\
S_2(x), & x\in[0.060;0.065],\\
S_3(x), & x\in[0.065;0.069],\\
S_4(x), & x\in[0.069;0.075],\\
S_5(x), & x\in[0.075;0.085],\\
S_6(x), & x\in[0.085;0.090],\\
S_7(x), & x\in[0.090;0.096].
\end{cases}
\]

где
\[
\begin{aligned}
S_0(x)={}&0.2079+0.0988(x-0.050)-0.5609(x-0.050)^2+582.3924(x-0.050)^3,\\
S_1(x)={}&0.2081+0.1035(x-0.052)+2.9335(x-0.052)^2-226.6399(x-0.052)^3,\\
S_2(x)={}&0.2090+0.1070(x-0.060)-2.5059(x-0.060)^2+222.8726(x-0.060)^3,\\
S_3(x)={}&0.2095+0.0986(x-0.065)+0.8372(x-0.065)^2-122.6772(x-0.065)^3,\\
S_4(x)={}&0.2099+0.0994(x-0.069)-0.6349(x-0.069)^2+121.8499(x-0.069)^3,\\
S_5(x)={}&0.2105+0.1050(x-0.075)+1.5584(x-0.075)^2-105.4709(x-0.075)^3,\\
S_6(x)={}&0.2116+0.1045(x-0.085)-1.6058(x-0.085)^2+141.5796(x-0.085)^3,\\
S_7(x)={}&0.2121+0.0991(x-0.090)+0.5179(x-0.090)^2-59.9346(x-0.090)^3.
\end{aligned}
\]

Проверим условия, налагаемые на сплайн.

\subsubsection*{1) Условия интерполяции}

Каждый многочлен $S_i(x)$ построен так, что
\[
S_i(x_i)=y_i, \quad S_i(x_{i+1})=y_{i+1}.
\]

Например,
\[
S_0(0.050)=0.2079, \quad S_0(0.052)=0.2081,
\]
\[
S_7(0.090)=0.2121, \quad S_7(0.096)=0.2127.
\]

Следовательно, условия интерполяции выполнены.

\subsubsection*{2) Непрерывность первой производной}

Производная частичного многочлена равна
\[
S_i'(x)=b_i+2c_i(x-x_i)+3d_i(x-x_i)^2.
\]

По построению
\[
S_i'(x_i)=m_i, \quad S_i'(x_{i+1})=m_{i+1}.
\]

Поэтому в каждом внутреннем узле выполняется условие
\[
S_{i-1}'(x_i)=S_i'(x_i)=m_i.
\]

Численно получаем
\[
\begin{aligned}
&S_0'(0.052)=S_1'(0.052)=0.1035,\\
&S_1'(0.060)=S_2'(0.060)=0.1070,\\
&S_2'(0.065)=S_3'(0.065)=0.0986,\\
&S_3'(0.069)=S_4'(0.069)=0.0994,\\
&S_4'(0.075)=S_5'(0.075)=0.1050,\\
&S_5'(0.085)=S_6'(0.085)=0.1045,\\
&S_6'(0.090)=S_7'(0.090)=0.0991.
\end{aligned}
\]

Следовательно, первая производная сплайна непрерывна.

\subsubsection*{3) Непрерывность второй производной}

Вторая производная частичного многочлена равна
\[
S_i''(x)=2c_i+6d_i(x-x_i).
\]

Проверим равенство вторых производных во внутренних узлах:
\[
\begin{aligned}
&S_0''(0.052)=S_1''(0.052)=5.8669,\\
&S_1''(0.060)=S_2''(0.060)=-5.0118,\\
&S_2''(0.065)=S_3''(0.065)=1.6744,\\
&S_3''(0.069)=S_4''(0.069)=-1.2699,\\
&S_4''(0.075)=S_5''(0.075)=3.1167,\\
&S_5''(0.085)=S_6''(0.085)=-3.2115,\\
&S_6''(0.090)=S_7''(0.090)=1.0358.
\end{aligned}
\]

Следовательно, вторая производная сплайна также непрерывна.

\subsubsection*{4) Проверка периодических граничных условий}

По условию периодичности должно выполняться равенство производных в крайних точках:
\[
S_0'(x_0)=S_7'(x_8).
\]

Действительно,
\[
S_0'(0.050)=0.0988, \quad S_7'(0.096)=0.0988.
\]

Также выполняется равенство вторых производных в крайних точках:
\[
S_0''(0.050)=S_7''(0.096)=-1.1218.
\]

Следовательно, построенный сплайн удовлетворяет периодическим граничным условиям.

\section{Вывод}

В результате выполнения лабораторной работы был построен интерполяционный кубический сплайн дефекта 1 для табличной функции из варианта 47. Были вычислены шаги таблицы, наклоны в узлах и коэффициенты всех частичных кубических многочленов. Также были проверены условия интерполяции, непрерывности первой и второй производных и периодические граничные условия.
